\documentclass[acmsmall,nonacm,screen]{acmart}

\usepackage{mathpartir}
\usepackage{booktabs}

\citestyle{acmauthoryear}
\settopmatter{printacmref=false}

\AtEndPreamble{\theoremstyle{acmdefinition}\newtheorem{remark}[theorem]{Remark}}

\newcommand{\kw}[1]{\mathsf{#1}}
\newcommand{\nm}[1]{\mathit{#1}}
\newcommand{\letin}[3]{\kw{let}\;#1 = #2\;\kw{in}\;#3}
\newcommand{\ite}[3]{\kw{if}\;#1\;\kw{then}\;#2\;\kw{else}\;#3}
\newcommand{\doo}[2]{\kw{do}\;\nm{#1}(#2)}
\newcommand{\onc}[2]{\kw{once}\;\nm{#1}(#2)}
\newcommand{\once}[5]{\kw{once}\;#1(#2)\;\kw{as}\;#3\;\kw{then}\;#4\;\kw{else}\;#5}
\newcommand{\raw}[2]{\kw{raw}\;\nm{#1}(#2)}
\newcommand{\gate}[3]{\kw{gate}\;#1\;\kw{then}\;#2\;\kw{else}\;#3}
\newcommand{\iter}[4]{\kw{iter}\;#1\;\kw{from}\;#2\;\kw{as}\;#3.\,#4}
\newcommand{\res}[1]{\kw{res}(#1)}
\newcommand{\lost}{\kw{lost}}
\newcommand{\ver}[1]{\kw{ver}(#1)}
\newcommand{\D}{\mathsf{D}}
\newcommand{\Vl}{\mathsf{V}}
\newcommand{\sem}[1]{[\![#1]\!]}
\newcommand{\evto}{\Downarrow}
\newcommand{\rqt}[1]{#1^{\circ}}
\newcommand{\strip}[2]{\lfloor #2 \rfloor_{#1}}
\newcommand{\rd}{\kw{read}}
\newcommand{\idem}{\kw{idem}}
\newcommand{\onceclass}{\kw{once}}

% A display for programs. It measures its content and warns when the content is wider than the line.
\newsavebox{\progbox}
\newenvironment{prog}
  {\begin{lrbox}{\progbox}$\displaystyle\begin{array}{@{}l@{\qquad}l@{}}}
  {\end{array}$\end{lrbox}%
   \ifdim\wd\progbox>\linewidth
     \PackageWarning{paper}{Overfull program display by \the\dimexpr\wd\progbox-\linewidth\relax}%
   \fi
   \[\usebox{\progbox}\]\ignorespacesafterend}

% written by journal.py; do not edit
\newcommand{\exAstray}{1,263}
\newcommand{\exBadPercent}{5}
\newcommand{\exChainCalls}{22.22}
\newcommand{\exCrashes}{23,709}
\newcommand{\exJournala}{10.1}
\newcommand{\exJournalb}{101.0}
\newcommand{\exJournalc}{1,010}
\newcommand{\exLate}{446}
\newcommand{\exLatePercent}{54}
\newcommand{\exLost}{3,360}
\newcommand{\exLucky}{1,592}
\newcommand{\exMarks}{24,559}
\newcommand{\exMedianNodes}{74}
\newcommand{\exPrograms}{4,000}
\newcommand{\exQuiet}{1,432}
\newcommand{\exQuietPercent}{72}
\newcommand{\exRefLost}{656}
\newcommand{\exReferences}{2,000}
\newcommand{\exRejected}{2,000}
\newcommand{\exRejectedRuns}{40,000}
\newcommand{\exRepeats}{3,823}
\newcommand{\exRequests}{74,714}
\newcommand{\exRules}{10}
\newcommand{\exRuns}{20,000}
\newcommand{\exScratcha}{10.6}
\newcommand{\exScratchb}{173.2}
\newcommand{\exScratchc}{2,316,257}
\newcommand{\exSlowest}{48,002}
\newcommand{\exStepLateReq}{4}
\newcommand{\exStepLenient}{8}
\newcommand{\exStepNoEpoch}{2}
\newcommand{\exStepNoLeftover}{392}
\newcommand{\exStepOnceRetry}{2}
\newcommand{\exSteps}{5}
\newcommand{\exStuck}{823}
\newcommand{\exTryDoArg}{107}
\newcommand{\exTryDoClass}{62}
\newcommand{\exTryDoPc}{624}
\newcommand{\exTryIfPc}{2,637}
\newcommand{\exTryIfResult}{48,002}
\newcommand{\exTryIter}{137}
\newcommand{\exTryOnceArg}{349}
\newcommand{\exTryOncePc}{21,051}
\newcommand{\exTryRawClass}{14}
\newcommand{\exTryRawLabel}{86}


\begin{document}

\title{A Crash Is a Lost Reply: Typed Journals for Agents That Restart}

\author{Claude-Fable-5.2}
\authornote{An AI model. The proofs are on paper and have not been checked by a machine. A
script implements the calculus and tests the theorems on random programs with random crashes.
At the time of writing nobody has reviewed the paper.}

\begin{abstract}
An agent that works for hours calls a model hundreds of times and acts on the world in
between, and somewhere on the way its process dies: a machine is reclaimed, or a new version
is deployed. Durable execution keeps such a run alive. Every request is written to a journal,
and a new process runs the program again from its start and takes the recorded answers from
the journal instead of asking again. This works only if the program makes the same requests
when it runs again, and it leaves open what becomes of the one request that was under way. We
give a calculus for such programs. A request is a $\kw{do}$, for an action that may be
repeated, or a $\kw{once}$, for an action that may not; a $\kw{once}$ has a second branch for
the case that its reply is lost. A type system with two labels, durable and volatile, keeps
what the journal does not hold, such as the clock, away from what it holds. We prove four
statements. A well-typed program never gets stuck in replay, whatever the crashes and the
versions. A run with crashes that ends is, as far as the world can tell, a run without
crashes in which a $\kw{once}$ may lose its reply and a $\kw{do}$ may be made several times in
a row; and every such run without crashes arises from some crashes. Each crash costs at most
one repeated request or one lost reply. And versions of a program that differ by gates can
replace one another in the middle of a run. So a program that handles a lost reply, as every
program that talks over a network must, handles a crash. A script implements the calculus and
tests the theorems on \exPrograms{} random programs with \exCrashes{} crashes; it refutes
\exRules{} weakened typing rules and \exSteps{} weakened steps of the protocol. Without a
journal, a task of a thousand calls with a one percent chance of a crash per call makes
\exScratchc{} calls on average; with one, \exJournalc{}.
\end{abstract}

\keywords{durable execution, replay, crash recovery, idempotence, information flow, agents}

\authorsaddresses{}

\maketitle

\section{Introduction}

An agent built on a language model is a loop. It calls the model, does what the model asks
for, and calls the model again with the result. A task of any size takes hundreds of turns
and runs for hours~\citep{tpp6}. A process does not reliably live that long. The machine is
reclaimed, memory runs out, a new version is deployed. If the agent then starts again from
nothing, it pays for every call a second time, and it does everything a second time: the
message is sent twice, and so is the money.

The remedy is known from workflow engines under the name of durable
execution~\citep{burckhardt2021,temporal-workflow}. The program writes every request and
every answer to a \emph{journal} that outlives the process. A new process runs the program
again from its start. Where the program makes a request that the journal already answers, it
takes the answer from the journal and nothing is sent. When the journal is used up, the
program is where its predecessor died, with the same values in its variables, and it goes on.

Replay puts three obligations on a program, and an agent breaks each of them easily.

\emph{The same requests.} The journal fits only if the program, run again, makes the requests
that it made before. A program that reads the clock, draws a random number, or calls a model
without the journal may take another path the second time. The journal then does not fit, and
the engine can only stop the run. Such a program passes every test in which nothing crashes.
It fails in production, on the night a machine is lost.

\emph{The request under way.} When a process dies between a request and the record of its
answer, nobody knows whether the action happened. Repeating it is right for a query and wrong
for a payment. Not repeating it leaves the program without an answer.

\emph{New versions.} A deployment is a crash followed by a restart with another program. The
prompt of an agent is part of that program. Whoever improves a prompt changes a request, and
the journals of all runs that are under way stop fitting.

This paper gives a calculus and a type system for these obligations, and it proves what a
program that meets them gets in return.

\paragraph{The calculus.} There are two forms of request. $\doo{a}{t}$ is for an action
that may be performed again: a call of the model, a query, a write that is idempotent. If its
answer was not recorded, the new process performs it again. The form
$\once{a}{t}{x}{e_1}{e_2}$ is for an action that must not be performed again, such as a
payment. If its answer was not recorded, the new process does not perform it; it records that
the reply is \emph{lost} and continues with $e_2$, where the program has to find out what
happened or ask. A third form, $\raw{a}{t}$, performs an action outside the journal. And
$\gate{v}{e_1}{e_2}$ runs $e_1$ in the parts of a run that began under version $v$ or later,
and $e_2$ in the older parts.

\paragraph{The types.} A value is \emph{durable} if the journal determines it and
\emph{volatile} otherwise. The result of a $\kw{raw}$ action is volatile, and so is everything
computed from it. The type system keeps volatile values out of the arguments of requests and
out of the tests that decide whether a request is made. It is a system of information flow
with two labels~\citep{volpano1996}, and what it protects is the journal.

\paragraph{Results.} We compare the machine, with its journal and its crashes, to a
\emph{reference semantics} that has neither: a program runs once, from start to end. The one
unusual rule of the reference semantics is that a $\kw{once}$ may lose its reply, as any
request over a network may.
\begin{itemize}
\item \emph{Replay never gets stuck} (Theorem~\ref{thm:stuck}). For a well-typed program the
  journal always fits, whatever the crashes, the replies, and the versions.
\item \emph{A crash is a lost reply} (Theorems~\ref{thm:sound} and~\ref{thm:complete}). A run
  with crashes that ends has the value and the world trace of a reference execution, in which
  a $\kw{do}$ may be performed several times in a row. Conversely, every reference execution
  is a run with suitable crashes. The first half holds for every program, typed or not: a
  journal that is checked strictly never lets a run end wrongly. It lets it get stuck.
\item \emph{A crash costs at most one} (Theorem~\ref{thm:cost}). The repeated performances and
  the lost replies of a run together number at most its crashes. No rule for the request
  under way can do better for an action that is not idempotent
  (Proposition~\ref{prop:exactly}).
\item \emph{Versions} (Theorem~\ref{thm:versions}). If all changes of a program are behind
  gates, the old texts are the new text with a lower version, and a run may change from one to
  the next at any crash. A change that is not behind a gate gets runs stuck, and so does a
  deployment that is taken back.
\end{itemize}
Together these say that a crash adds nothing to what a careful program has to handle anyway.
A program that deals with a lost reply at every $\kw{once}$, and that keeps the clock away
from its requests, cannot tell a crash from a slow network.

\paragraph{Tests.} The proofs are on paper. An artifact implements the checker, the machine,
and the reference semantics, and tests each theorem on random programs with random crashes
and deployments (Section~\ref{sec:tests}). It weakens the type system in \exRules{} ways and
the protocol of the journal in \exSteps{} ways, and finds for each a run that the theorems
exclude. It also shows why these defects are found late: of the programs that the checker
rejects, \exQuietPercent{} percent show nothing wrong in twenty runs with crashes.

\section{One Program, Three Runs}
\label{sec:overview}

\subsection{A Reply Is Lost}
\label{sec:booking}

The agent of this section books a journey. The model picks a fare, the fare is paid, and the
receipt is filed.
\begin{prog}
\letin{f}{\doo{model}{0}}{} & \text{the model picks a fare}\\
\kw{once}\;\nm{pay}(f)\;\kw{as}\;r\;\kw{then}\;\doo{put}{r} & \text{pay, and file the receipt}\\
\phantom{\kw{once}\;\nm{pay}(f)\;\kw{as}\;r}\;\kw{else}\;\letin{s}{\doo{look}{f}}{\doo{put}{s}} & \text{the reply is lost: look it up, then file}
\end{prog}
The actions have classes. $\nm{model}$ and $\nm{look}$ are of class $\rd$: they change
nothing. $\nm{put}$ is of class $\idem$: filing the same receipt twice in a row is the same
as filing it once. $\nm{pay}$ is of class $\onceclass$.

Suppose the process dies while the payment is under way. The journal then reads
\[
\ver{1} \cdot \doo{model}{0} \cdot \res{2} \cdot \onc{pay}{2}.
\]
The first entry marks the version of the program. The model was asked and answered $2$. The
payment was requested, and that is the last thing the journal knows. The request was written
before the payment was made, so the payment may have happened or not. In this run it
happened, and the reply went down with the process.

A new process runs the program from its start. At $\doo{model}{0}$ it finds the request and
the answer $2$ in the journal, so the model is not called, and $f$ is $2$ as before. At
$\onc{pay}{f}$ it finds the request without an answer. It does not pay again. It writes
$\lost$ and takes the second branch, which looks up the payment and files what it finds. At
the end the journal reads
\[
\ver{1} \cdot \doo{model}{0} \cdot \res{2} \cdot \onc{pay}{2} \cdot \lost \cdot \doo{look}{2}
  \cdot \res{1} \cdot \doo{put}{1} \cdot \res{0},
\]
and the world has seen four actions: $\nm{model}(0)$, $\nm{pay}(2)$, $\nm{look}(2)$, and
$\nm{put}(1)$. Each happened once.

Nothing in this trace of the world says that a process died. It is the trace of a run without
any crash in which the reply to the payment did not arrive. A program that pays over a network
has to be ready for that, and the second branch of $\kw{once}$ makes the type checker insist
on it. Theorem~\ref{thm:sound} says that this is general: for the world, a crash is a lost
reply, and nothing else.

Had the process died one step later, after the reply $\res{3}$ was written, the new process
would have taken the first branch with $r = 3$. Had it died while the model was being called,
the new process would have called the model again, which is harmless and costs one call.

\subsection{What Must Not Reach the Journal}
\label{sec:volatile}

\begin{table}[t]
\caption{Programs of Section~\ref{sec:overview} and what the checker of the artifact says. A
verdict other than ``accepted'' is the reason for which the checker rejects the program.}
\label{tab:verdicts}
\small
\begin{tabular}{@{}ll@{}}
\toprule
Program & Verdict \\
\midrule
the booking program of \S\ref{sec:booking} & accepted \\
$\letin{f}{\doo{model}{0}}{\doo{pay}{f}}$ & an action that must not be repeated \\
$\letin{t}{\raw{clock}{0}}{\doo{put}{t}}$ & a request with a volatile argument \\
$\letin{t}{\raw{clock}{0}}{\ite{t \le 1}{\doo{put}{1}}{0}}$ & a request under volatile control \\
$\letin{t}{\doo{clock}{0}}{\ite{t \le 1}{\doo{put}{t}}{0}}$ & accepted \\
$\letin{t}{\raw{clock}{0}}{\letin{u}{\raw{look}{t}}{\doo{put}{1}}}$ & accepted \\
$\raw{mail}{1}$ & an effect outside the journal \\
the booking program with the gate of \S\ref{sec:prompt} & accepted \\
\bottomrule
\end{tabular}
\end{table}

A $\kw{raw}$ action goes around the journal. It is performed whenever the program comes to
it, in the first process and in every later one, and its result may differ each time. The
third program in Table~\ref{tab:verdicts} reads the clock in this way and files the time.
Suppose the process dies after the request $\doo{put}{2}$ is written. The next process reads
another time, say $3$, and comes to the request $\doo{put}{3}$ where the journal has
$\doo{put}{2}$. What should it do? If it trusts the journal, the program goes on believing
that it filed $3$ while the world holds $2$. So the machine compares, finds that the journal
does not fit, and stops: the run is \emph{stuck}. The fourth program is stuck in the same way
when the clock decides whether a request is made at all.

The checker rejects both programs. The value $t$ is volatile, and a volatile value may be
neither the argument of a request nor the test above one. The fifth program reads the clock
through the journal, with $\kw{do}$. The time is then recorded once and is the same in every
process. The sixth program shows that volatile values are not forbidden: they may flow into
other $\kw{raw}$ actions, such as a log, and into the result. They may only not flow into the
journal.

The second and the seventh program break the other rule, which concerns classes. A $\kw{do}$
may be performed twice, so it must not pay. A $\kw{raw}$ action is performed by every process,
so it must not have an effect at all.

\subsection{A Prompt Is Code}
\label{sec:prompt}

The argument of $\nm{model}$ stands for the prompt. Suppose the author improves it, so that
the first line of the program becomes $\letin{f}{\doo{model}{1}}{}$, and deploys the new
text while a run is under way whose journal begins with
$\ver{1} \cdot \doo{model}{0} \cdot \res{2}$. The next process of that run asks
$\doo{model}{1}$ where the journal has $\doo{model}{0}$, and the run is stuck. An edit of a
prompt is a change of the program, and the journal notices.

The remedy is to keep both versions in the text, behind a gate:
\begin{prog}
\letin{f}{\gate{2}{\doo{model}{1}}{\doo{model}{0}}}{\ldots}
\end{prog}
A gate asks under which version the present part of the run began. The journal knows: it
holds a mark $\ver{v}$ wherever a process with a newer version first went beyond the end of
the journal. In the run above, the new process comes to the gate with the cursor at the
start, where the mark $\ver{1}$ stands. So it takes the old branch, and its request fits the
journal. When it reaches the end of the journal it writes the mark $\ver{2}$, and every gate
after that point takes its new branch. A run that starts after the deployment takes the new
branch everywhere. Theorem~\ref{thm:versions} says that this is safe in general, and that the
old text is the new text run with the old version. The versions must not go down: if the old
text is deployed again after a run has passed the mark $\ver{2}$, the old text has nothing to
say about the requests behind the mark, and the run is stuck.

\section{The Calculus}
\label{sec:calculus}

\subsection{Syntax}

\begin{figure}[t]
\[
\begin{array}{@{}l@{\quad}r@{\;}l@{}}
\text{terms} & t ::= & x \mid n \mid t + t \\
\text{tests} & b ::= & t = t \mid t \le t \\
\text{expressions} & e ::= & t \mid \letin{x}{e}{e} \mid \ite{b}{e}{e} \mid \iter{n}{t}{x}{e} \\
 & \mid & \kw{do}\;a(t) \mid \once{a}{t}{x}{e}{e} \mid \kw{raw}\;a(t) \mid \gate{v}{e}{e} \\[1ex]
\text{entries} & & \ver{v} \mid \kw{do}\;a(c) \mid \kw{once}\;a(c) \mid \res{r} \mid \lost \\
\text{events of the world} & & a(c){:}r \mid \kw{raw}\;a(c){:}r
\end{array}
\]
\Description{The grammar of terms, tests, expressions, journal entries, and events.}
\caption{Syntax. The letter $a$ ranges over actions, each of which has a class, $\rd$,
$\idem$, or $\onceclass$; $n$, $c$, and $r$ range over numbers, $x$ over variables, and $v$
over versions, which are positive numbers.}
\label{fig:syntax}
\end{figure}

Figure~\ref{fig:syntax} gives the syntax. Values are numbers. An action takes a number and
replies with a number; the number stands for whatever is sent and received. Every action has a
\emph{class}. An action of class $\rd$ changes nothing in the world. An action of class
$\idem$ has an effect, and performing it several times in a row, with the same argument, has
the effect of performing it once. An action of class $\onceclass$ has an effect that counts
every time. The classes are claims about the world, and the calculus takes them as given.

Loops are bounded and there are no functions, so every program ends. A loop
$\iter{n}{t}{x}{e}$ has a state $x$. The state starts as the value of $t$, and each of the $n$
turns evaluates $e$ and makes its value the new state.

\subsection{The Machine}
\label{sec:machine}

A \emph{journal} is a sequence of entries: marks $\ver{v}$, requests $\kw{do}\;a(c)$ and
$\kw{once}\;a(c)$, and answers $\res{r}$ and $\lost$. A \emph{world trace} is a sequence of
events: $a(c){:}r$ for an action that was performed for a request, and $\kw{raw}\;a(c){:}r$
for one that was performed outside the journal. The journal $J$ and the world trace $W$
outlive a process. Everything else is lost when a process dies.

An \emph{attempt} is one process. It is given the journal, the world trace, and the version
$\hat v$ under which it is deployed. It evaluates the program from its start, with a
\emph{cursor} $i$ into the journal, at first $0$, and a \emph{current version} $u$, at first
$0$. Variables, tests, and loops are evaluated as usual. The other forms use the journal.

\begin{table}[t]
\caption{What a request $\kappa\;a(c)$ does, where $\kappa$ is $\kw{do}$ or $\kw{once}$,
after the version has been brought up to date. ``This request'' is the entry $\kappa\;a(c)$.
The steps are made from left to right, and the process may die before any of them.}
\label{tab:protocol}
\small
\begin{tabular}{@{}llll@{}}
\toprule
 & At the cursor & Steps & Answer \\
\midrule
1 & the end of the journal & write this request; perform $a(c)$, reply $r$; write $\res{r}$ & $r$ \\
2 & this request, then $\res{r}$ & none & $r$ \\
3 & this request, then $\lost$ & none & lost \\
4 & this request as the last entry, $\kappa = \kw{do}$ & perform $a(c)$, reply $r$; write $\res{r}$ & $r$ \\
5 & this request as the last entry, $\kappa = \kw{once}$ & write $\lost$ & lost \\
6 & anything else & none & stuck \\
\bottomrule
\end{tabular}
\end{table}

\paragraph{Versions.} Before a request and at a gate, the attempt brings its version up to
date. While the entry at the cursor is a mark $\ver{v}$, it sets $u$ to $v$ and moves the
cursor past the mark. If the cursor is then at the end of the journal and $u < \hat v$, it
writes the mark $\ver{\hat v}$ and sets $u$ to $\hat v$. So a mark stands where a newer
version first went beyond what older versions had recorded.

\paragraph{Requests.} Table~\ref{tab:protocol} says what a request does. At the end of the
journal the attempt is \emph{live}: it writes the request, performs the action, and writes the
reply (row 1). The request is written first, so that a process which dies while the action is
under way leaves a trace of it. Before the end of the journal the attempt \emph{replays}: if
the journal holds this request with an answer, the answer is taken and nothing is performed
(rows 2 and 3). If the journal holds this request as its last entry, the attempt that wrote
it died before it could write an answer. A $\kw{do}$ is then performed, perhaps for the
second time (row 4). A $\kw{once}$ is not: its reply is recorded as lost (row 5). If the
entry at the cursor is anything but this request, the journal does not fit the program, and
the run is stuck (row 6). The value of $\kw{do}\;a(t)$ is the reply. The expression
$\once{a}{t}{x}{e_1}{e_2}$ goes on with $e_1$, where $x$ is the reply, or with $e_2$ if the
reply is lost. Every write moves the cursor to the new end of the journal, and every
performed action adds the event $a(c){:}r$ to the world trace.

\paragraph{The other forms.} $\gate{v}{e_1}{e_2}$ brings the version up to date and goes on
with $e_1$ if $u \ge v$ and with $e_2$ otherwise. $\kw{raw}\;a(t)$ performs the action and
adds $\kw{raw}\;a(c){:}r$ to the world trace; its value is the reply. When the program has
its value, the attempt moves the cursor past any marks. If the cursor is then not at the end
of the journal, the run is stuck: the journal holds requests that this attempt did not make.

\paragraph{Runs.} An attempt ends in one of three ways: with a value, stuck, or because the
process dies, which can happen at any moment and leaves $J$ and $W$ as they are. A \emph{run}
is a sequence of attempts on the same program. The first starts with an empty journal and an
empty world trace. Each of the others starts with the $J$ and $W$ that its predecessor left,
under a version of its own. All attempts but the last die, and the last one ends with a value
or stuck. The replies of the world are arbitrary, so a statement about all runs is a
statement about every world, every model, and every way in which processes die.

A journal of a run always has a simple shape, which we call \emph{well formed}: marks, and
requests each followed by its answer, where $\lost$ answers only a $\kw{once}$; the last entry
may be a request without an answer, which we call \emph{open}.

\subsection{The Reference Semantics}
\label{sec:reference}

\begin{figure}[t]
\begin{mathpar}
\inferrule*[right=Term]{ }{\rho; k \vdash t \evto \sem{t}\rho; k; \epsilon}
\and
\inferrule*[right=Raw]{c = \sem{t}\rho}{\rho; k \vdash \kw{raw}\;a(t) \evto r; k; \kw{raw}\;a(c){:}r}
\and
\inferrule*[right=Do]{c = \sem{t}\rho}{\rho; k \vdash \kw{do}\;a(t) \evto r; k + 1; \kw{do}\;a(c){:}r}
\and
\inferrule*[right=Let]{\rho; k \vdash e_1 \evto v_1; k_1; T_1 \\ \rho[x \mapsto v_1]; k_1 \vdash e_2 \evto v_2; k_2; T_2}
  {\rho; k \vdash \letin{x}{e_1}{e_2} \evto v_2; k_2; T_1 T_2}
\and
\inferrule*[right=Once-Res]{c = \sem{t}\rho \\ \rho[x \mapsto r]; k + 1 \vdash e_1 \evto v; k'; T}
  {\rho; k \vdash \once{a}{t}{x}{e_1}{e_2} \evto v; k'; \kw{once}\;a(c){:}r \cdot T}
\and
\inferrule*[right=Once-Lost]{c = \sem{t}\rho \\ \rho; k + 1 \vdash e_2 \evto v; k'; T}
  {\rho; k \vdash \once{a}{t}{x}{e_1}{e_2} \evto v; k'; \kw{lost}\;a(c) \cdot T}
\and
\inferrule*[right=Gate-New]{\kappa(k) \ge v \\ \rho; k \vdash e_1 \evto v'; k'; T}
  {\rho; k \vdash \gate{v}{e_1}{e_2} \evto v'; k'; T}
\and
\inferrule*[right=Gate-Old]{\kappa(k) < v \\ \rho; k \vdash e_2 \evto v'; k'; T}
  {\rho; k \vdash \gate{v}{e_1}{e_2} \evto v'; k'; T}
\end{mathpar}
\Description{Eight rules of a big-step semantics with a counter of requests and a trace.}
\caption{The reference semantics, for a fixed schedule $\kappa$ of versions. The rules for
conditionals and loops are the usual ones; like \textsc{Let}, they pass the counter along and
join the traces. The reply $r$ in \textsc{Raw}, \textsc{Do}, and \textsc{Once-Res} is any
number.}
\label{fig:reference}
\end{figure}

The reference semantics says what a program does when nothing crashes. It has no journal and
no cursor. The judgment $\rho; k \vdash e \evto v; k'; T$ of Figure~\ref{fig:reference} says
that $e$, in the environment $\rho$ and after $k$ requests, has the value $v$ after $k'$
requests, with the trace $T$. A trace is a sequence of events $\kw{do}\;a(c){:}r$,
$\kw{once}\;a(c){:}r$, $\kw{lost}\;a(c)$, and $\kw{raw}\;a(c){:}r$. We write $\rqt{T}$ for $T$
without its $\kw{raw}$ events and call it the \emph{request trace}.

Two rules are not the usual ones. \textsc{Once-Lost} lets a $\kw{once}$ go without a reply.
The rule can always be applied, and it does not say whether the action happened. It models
what any program that sends a request over a network must expect. The two rules for gates
read the version from a \emph{schedule} $\kappa$, a function from numbers to versions that
never goes down: $\kappa(k)$ is the version that is in force after $k$ requests. A schedule
that is constant describes a run under one version.

\begin{definition}[Traces of the world]
\label{def:trace}
A world trace $W$ \emph{belongs to} a trace $T$ if the events of $W$ that are not $\kw{raw}$
events are, in order, the following: for every event $\kw{do}\;a(c){:}r$ of $T$, one or more
events $a(c){:}r'$, of which the last has the reply $r$; for every event
$\kw{once}\;a(c){:}r$, the one event $a(c){:}r$; and for every event $\kw{lost}\;a(c)$, at
most one event $a(c){:}r'$.
\end{definition}

So a $\kw{do}$ may be performed several times, but only in a row, with nothing else between
the performances except $\kw{raw}$ actions. For an action of class $\rd$ or $\idem$ that is
the same as performing it once, by what the classes mean.

\section{Types}
\label{sec:types}

\begin{figure}[t]
\begin{mathpar}
\inferrule*[right=Term]{\Gamma \vdash t : \ell}{\Gamma; pc \vdash t : \ell}
\and
\inferrule*[right=Raw]{\text{$a$ has class $\rd$}}{\Gamma; pc \vdash \kw{raw}\;a(t) : \Vl}
\and
\inferrule*[right=Do]{\Gamma \vdash t : \D \\ pc = \D \\ \text{$a$ has class $\rd$ or $\idem$}}
  {\Gamma; pc \vdash \kw{do}\;a(t) : \D}
\and
\inferrule*[right=Once]{\Gamma \vdash t : \D \\ pc = \D \\ \Gamma, x{:}\D; pc \vdash e_1 : \ell_1 \\ \Gamma; pc \vdash e_2 : \ell_2}
  {\Gamma; pc \vdash \once{a}{t}{x}{e_1}{e_2} : \ell_1 \sqcup \ell_2}
\and
\inferrule*[right=Let]{\Gamma; pc \vdash e_1 : \ell_1 \\ \Gamma, x{:}\ell_1; pc \vdash e_2 : \ell_2}
  {\Gamma; pc \vdash \letin{x}{e_1}{e_2} : \ell_2}
\and
\inferrule*[right=If]{\Gamma \vdash b : \ell \\ \Gamma; pc \sqcup \ell \vdash e_1 : \ell_1 \\ \Gamma; pc \sqcup \ell \vdash e_2 : \ell_2}
  {\Gamma; pc \vdash \ite{b}{e_1}{e_2} : \ell \sqcup \ell_1 \sqcup \ell_2}
\and
\inferrule*[right=Gate]{\Gamma; pc \vdash e_1 : \ell_1 \\ \Gamma; pc \vdash e_2 : \ell_2}
  {\Gamma; pc \vdash \gate{v}{e_1}{e_2} : \ell_1 \sqcup \ell_2}
\and
\inferrule*[right=Iter]{\Gamma \vdash t : \ell_0 \\ \Gamma, x{:}\ell; pc \vdash e : \ell' \\ \ell_0 \sqcup \ell' \sqsubseteq \ell}
  {\Gamma; pc \vdash \iter{n}{t}{x}{e} : \ell}
\end{mathpar}
\Description{Eight typing rules with labels D and V and a program counter label.}
\caption{Typing. Labels are $\D$, durable, and $\Vl$, volatile, with
$\D \sqsubseteq \Vl$. The judgment $\Gamma \vdash t : \ell$ says that $\ell$ is the join of the
labels of the variables of the term $t$, which is $\D$ if there are none; the label of a test
is the join of the labels of its two terms.}
\label{fig:typing}
\end{figure}

There are two labels, $\D$ for durable and $\Vl$ for volatile, with $\D$ below $\Vl$. The
judgment $\Gamma; pc \vdash e : \ell$ of Figure~\ref{fig:typing} gives the expression $e$ the
label $\ell$, where $\Gamma$ gives a label to every variable and $pc$ is the label of the
tests that decide whether $e$ is evaluated. A program is \emph{well typed} if it has a label
with no variables and $pc = \D$.

The rules are those of a type system for information flow~\citep{volpano1996,sabelfeld2003},
with the journal in the place of the public output. The reply to a request is durable: the
journal holds it, and every process reads the same. The reply to a $\kw{raw}$ action is
volatile. \textsc{If} raises $pc$ to the label of its test, and \textsc{Do} and \textsc{Once}
ask for $pc = \D$ and a durable argument. So what a program asks, and whether it asks, depend
on durable values only.

A gate does not raise $pc$. The version that it reads comes from the marks of the journal,
and it is as durable as a reply. A gate may also stand under a volatile test. It then writes
a mark in one process and perhaps not in another, and Section~\ref{sec:results} shows that
this does no harm, because a mark is read by position and never compared.

The classes enter in two rules. \textsc{Do} excludes actions of class $\onceclass$, since row
4 of Table~\ref{tab:protocol} may perform a $\kw{do}$ twice. \textsc{Raw} admits only class
$\rd$, since every process of a run performs the $\kw{raw}$ actions on its path. A
$\kw{once}$ may be used for any class; for an action that is idempotent it is merely more
careful than necessary.

The checker of the artifact computes the least label of an expression, and for a loop the
least label of its state. It has fewer than fifty lines.

\section{Results}
\label{sec:results}

\subsection{Requests Are Durable}

An execution of the reference semantics can be cut short. A \emph{stopped} execution is a
derivation that ends at some point before its value; its trace is the trace up to that
point. If it stops while a request waits for its answer, the trace ends with an \emph{open}
event $\kw{do}\;a(c){:}{?}$ or $\kw{once}\;a(c){:}{?}$. Two events are \emph{answers to the
same request} if they agree in everything but the reply, where $\kw{lost}\;a(c)$ counts as an
answer to $\kw{once}\;a(c)$. An open event counts as the beginning of every answer to its
request: a trace that ends with it is a prefix of the traces that go on with an answer.

\begin{lemma}[Requests are durable]
\label{lem:durable}
Let $\Gamma; pc \vdash e : \ell$, and let $\rho_1$ and $\rho_2$ agree on the variables that
$\Gamma$ labels $\D$. Take two executions of $e$, complete or stopped, in $\rho_1$ and in
$\rho_2$, with the same schedule and the same counter $k$, and let $T_1$ and $T_2$ be their
traces.
\begin{enumerate}
\item Either one of $\rqt{T_1}$ and $\rqt{T_2}$ is a prefix of the other, or they have the forms
  $S \cdot \alpha_1 \cdot S_1$ and $S \cdot \alpha_2 \cdot S_2$ for two different answers
  $\alpha_1$ and $\alpha_2$ to the same request.
\item If the first execution is complete and $\rqt{T_1}$ is a prefix of $\rqt{T_2}$, then
  $\rqt{T_1} = \rqt{T_2}$. If the second is complete as well, the two end with the same counter,
  and with the same value if $\ell = \D$.
\item If $pc = \Vl$, then $\rqt{T_1}$ and $\rqt{T_2}$ are empty.
\end{enumerate}
\end{lemma}

\begin{proof}
(3) holds because \textsc{Do} and \textsc{Once} need $pc = \D$, and no rule lowers $pc$
towards a subexpression. For (1) and (2) we go by induction on $e$. We show the cases for
complete executions; a stopped execution stops inside a subexpression, where the induction
hypothesis applies, or between two, where its trace is a prefix.

\emph{Term and Raw.} There are no requests. A term with the label $\D$ has only variables on
which the environments agree. A $\kw{raw}$ action has the label $\Vl$.

\emph{Do.} The argument is durable, so it is the same number $c$ in both executions. The
request traces are $\kw{do}\;a(c){:}r_1$ and $\kw{do}\;a(c){:}r_2$. If the replies are the
same, so are the traces and the values. Otherwise they are two answers to the same request.

\emph{Once.} The request is the same, as for \textsc{Do}. If the answers differ, we are done.
If both are the reply $r$, the executions go on with $e_1$ in environments that agree on
durable variables, $x$ among them, from the same counter. If both are lost, they go on with
$e_2$. The induction hypothesis applies.

\emph{Gate.} Both executions read $\kappa(k)$ and take the same branch.

\emph{If.} If the test is durable, both executions take the same branch. If it is volatile,
the branches are typed with $pc = \Vl$ and make no requests by (3), and the label of the
conditional is $\Vl$.

\emph{Let.} By the induction hypothesis for $e_1$, the request traces of $e_1$ differ in an
answer, and then so do the whole traces, or they are equal. In the second case the counters
are equal, and the values are equal if the label of $e_1$ is $\D$. The environments for $e_2$
then agree on durable variables, and the induction hypothesis applies to $e_2$.

\emph{Iter.} By induction on the number of turns, each turn as in \textsc{Let}. The label
$\ell$ of the state is above the label of the body, so the states of the two executions agree
after every turn if $\ell = \D$.
\end{proof}

The lemma is a noninterference property: the requests of a program do not depend on its
volatile values. What is specific to journals is the form of the conclusion. Two executions
of a well-typed program can differ in their requests only \emph{after} they have received
different answers.

\subsection{Replay Never Gets Stuck}

\begin{definition}[Recording]
\label{def:records}
A journal $J$ \emph{records} an execution of $e$, complete or stopped, with the schedule
$\kappa$, if $J$ is well formed; the requests of $J$ with their answers are, in order, the
events of the request trace of the execution, where a request without an answer in $J$ is an
open event; and for every $k$ below the number of requests of $J$, $\kappa(k)$ is the last
mark of $J$ before its request number $k + 1$.
\end{definition}

\begin{theorem}[Replay never gets stuck]
\label{thm:stuck}
Let $e$ be well typed. Then no attempt of a run of $e$ is stuck, whatever the replies, the
versions, and the moments at which processes die. After every attempt, the journal records an
execution of $e$.
\end{theorem}

\begin{proof}
By induction on the number of attempts. The empty journal records the execution that is
stopped at its start. Let $J$ record an execution $X$, and let $B$ be the next attempt. As
long as it runs, $B$ is an execution of $e$ in the reference semantics, stopped where $B$
is: its answers are the ones that it reads or writes, its $\kw{raw}$ replies are the ones
that it receives, and its schedule gives, after $k$ requests, the version $u$ that it holds
once its version is up to date. While the cursor of $B$ is before the end of $J$, the answers
that $B$ has read are those of $X$, and the versions that it has held are those of the
schedule of $X$, by Definition~\ref{def:records}.

\emph{$B$ is not stuck at a request.} Let $B$ come to a request after $m$ requests, with the
cursor before the end of $J$. The request trace of $B$, stopped here, consists of $m$
answered requests and an open one. The $m$ answered requests are the first $m$ events of the
request trace of $X$, which has a further event, because $J$ has a further request. The
schedules of $B$ and of $X$ agree up to this request, and the lemma uses no more of them. By
Lemma~\ref{lem:durable}(1) the two events number $m + 1$ are answers to the same request. So
the entry at the cursor is the request of $B$, and row 6 of Table~\ref{tab:protocol} does not
apply for this reason. Nor does it apply because of the answer: if the entry is followed by
$\lost$, it is a $\kw{once}$, because $J$ is well formed.

\emph{$B$ is not stuck at its end.} If $B$ had its value with requests of $J$ left over, the
request trace of $B$, which is complete, would be a proper prefix of that of $X$. This
contradicts Lemma~\ref{lem:durable}(2).

\emph{The journal after $B$.} If $B$ dies before it writes anything, the journal is $J$ and
records $X$. Otherwise consider the first write of $B$. At that moment the cursor is at the
end of $J$, or $J$ ends with an open request that $B$ has just matched. All requests and
answers of $J$ have been matched by $B$, and all its marks have been read by $B$ at the
places where they stand. So $J$ records the execution of $B$ up to this point, and from here
on $B$ writes what it does: each mark where its version rises, each request before it is
made, each answer when it is known. Whenever $B$ dies or ends, the journal records the
execution of $B$, stopped at its last write or complete.
\end{proof}

The theorem needs no order among the versions. With one program text, a process of an older
version writes no mark and follows the marks that are there.

\subsection{A Crash Is a Lost Reply}

\begin{lemma}[The journal counts]
\label{lem:counts}
At every moment of a run, of any program, the events of $W$ that are not $\kw{raw}$ events
fall into consecutive groups, one for each request of $J$ in order, such that the group of a
request for $a(c)$ consists of events $a(c){:}r'$ and
\begin{itemize}
\item for a $\kw{do}$ with the answer $\res{r}$, the group is not empty and its last reply is
  $r$; for a $\kw{do}$ that is open, the group may be empty;
\item for a $\kw{once}$ with the answer $\res{r}$, the group is the one event $a(c){:}r$;
\item for a $\kw{once}$ that is open or answered by $\lost$, the group has at most one event.
\end{itemize}
\end{lemma}

\begin{proof}
By Table~\ref{tab:protocol}. An action is performed for a request in two places only. In row
1 the request has just been written at the end of the journal, and the reply is written next.
In row 4 the request is a $\kw{do}$ and the last entry of the journal. In both, the request is
the last one of the journal, so its events come after those of all other requests, and they
are consecutive. A $\kw{once}$ is performed in row 1 only, by the attempt that writes it.
\end{proof}

\begin{theorem}[A crash is a lost reply]
\label{thm:sound}
Let a run of a program $e$, well typed or not, end with the value $v$, the journal $J$, and
the world trace $W$. Then there is a complete execution of $e$ in the reference semantics
that has the value $v$, that is recorded by $J$, and to whose trace $W$ belongs.

If $e$ is well typed, then every run in which only finitely many processes die ends with a
value; in its world trace, the events of a request for an action of class $\onceclass$ are at
most one, and all $\kw{raw}$ events are of class $\rd$.
\end{theorem}

\begin{proof}
Let $B$ be the last attempt. It ends with $v$ and is not stuck, so it has matched every
request of the journal that it found, has read every mark, and has written the rest. As in
the proof of Theorem~\ref{thm:stuck}, $B$ is a complete execution in the reference semantics,
and $J$ records it; that part of the proof does not use the types. By
Lemma~\ref{lem:counts}, the events of $W$ fall into groups that are exactly what
Definition~\ref{def:trace} asks of the trace of this execution: an answer $\res{r}$ to a
$\kw{do}$ is the event $\kw{do}\;a(c){:}r$, and so on.

For a well-typed program, no attempt is stuck by Theorem~\ref{thm:stuck}, and every attempt
that does not die ends, because programs end. The rule \textsc{Do} excludes the class
$\onceclass$, so by Lemma~\ref{lem:counts} only a $\kw{once}$ performs such an action, at most
once for a request. The rule \textsc{Raw} admits the class $\rd$ only.
\end{proof}

The first half of the theorem is about every program. A journal that is compared strictly
never lets a run end with a wrong account of what happened: if the program has left the path
that the journal recorded, the run gets stuck. So the strict comparison alone gives safety,
at run time. What the types add is that the run does not get stuck. They are about progress.

The theorem is also tight.

\begin{theorem}[Every reference execution is a run]
\label{thm:complete}
Let $e$ have a complete execution in the reference semantics with the value $v$, the trace
$T$, and a schedule $\kappa$ with $\kappa(0) \ge 1$. Let $W$ be a world trace without
$\kw{raw}$ events that belongs to $T$ and has exactly one event for each $\kw{do}$ event of
$T$. Then there is a run of $e$ that ends with the value $v$, whose world trace is $W$ apart
from $\kw{raw}$ events, and whose journal records the execution.
\end{theorem}

\begin{proof}
We choose the replies, the versions, and the deaths. The world gives the replies of $T$, and
it gives every attempt the same replies to its $\kw{raw}$ actions, in the same order. Then
every attempt follows the path of the execution as far as it gets, since the reference
semantics is deterministic once the answers, the versions, and the $\kw{raw}$ replies are
fixed. The first attempt has the version $\kappa(0)$. A process dies in two situations. When
it has written a $\kw{once}$ whose answer in $T$ is lost, it dies before it performs the
action or just after, according to whether $W$ has an event for it; the next attempt writes
$\lost$ by row 5. And when it has made $k$ requests and $\kappa(k)$ is above its version, it
dies before its next step, and the next attempt has the version $\kappa(k)$ and writes its
mark at this place. Nothing is performed twice, so the world trace is $W$ apart from the
$\kw{raw}$ events of the attempts.
\end{proof}

So the runs of a program under crashes and the executions of the reference semantics are the
same, as far as values and the world are concerned, except that a $\kw{do}$ may happen more
than once in a row. Whoever wants to know what a program can do when processes die may forget
the journal and read the program with the rule \textsc{Once-Lost} in mind.

\subsection{What a Crash Costs}

\begin{theorem}[A crash costs at most one]
\label{thm:cost}
In every run, the number of performances of $\kw{do}$ requests beyond the first of each
request, plus the number of answers $\lost$, is at most the number of processes that died.
\end{theorem}

\begin{proof}
A request is open from the moment it is written until its answer is written, and at most one
request is open at a time. We charge to each request the processes that die while it is open.
Every attempt performs a given request at most once. If a $\kw{do}$ is performed by $n$
attempts, each of them but the last died after performing it and before writing the reply, so
$n - 1$ deaths are charged to it. If a $\kw{once}$ is answered by $\lost$, the answer was
written in row 5 by an attempt that found the request open, so at least one death is charged
to it. A $\kw{once}$ is not performed twice, and a $\kw{do}$ is not lost.
\end{proof}

For an agent, most requests are calls of the model. A crash wastes at most the one call that
was under way. It wastes nothing if the process died between two calls.

Could the machine do better than rows 4 and 5, and perform the action exactly once?

\begin{proposition}[No exactly once]
\label{prop:exactly}
Consider an action that can only be performed, and a rule that decides from the journal
whether the new process performs an open request for it. No such rule has the action
performed exactly once in every run.
\end{proposition}

\begin{proof}
Take two runs. In one the process dies just before it performs the action, in the other
just after. Both leave the same journal, which ends with the open request. A rule that
performs the action performs it a second time in the second run. A rule that does not leaves
it undone in the first.
\end{proof}

This is the old observation that a failure between an action and its record cannot be hidden
from the application~\citep{saltzer1984,elnozahy2002}. The two forms of request are the two
ways out. $\kw{do}$ performs again and is right if repetition is harmless. $\kw{once}$ does
not, and tells the program. A world that lets the program look, as $\nm{look}$ does in
Section~\ref{sec:booking}, or that accepts a key with each request and ignores a key that it
has seen, turns the second case into the first.

\subsection{Versions}

Let $\strip{c}{e}$ be the program $e$ in which every gate with a version above $c$ is replaced
by its second branch. We read $e$ as the newest text of a program, and $\strip{c}{e}$ as the
text that was deployed as version $c$. An author who puts every change behind a gate with a
new version number produces texts that are related in this way.

A mark of a journal is \emph{redundant} if no request stands between it and the next mark or
the end.

\begin{theorem}[Versions]
\label{thm:versions}
Let $c_1 \le c_2 \le \cdots$ be versions. Consider the runs in which attempt $j$ executes the
text $\strip{c_j}{e}$ under the version $c_j$, and the runs in which attempt $j$ executes $e$
under the version $c_j$. Both have the same values, the same world traces, and the same
journals up to redundant marks. If $e$ is well typed, then so is every $\strip{c}{e}$, and no
run of the first kind is stuck.
\end{theorem}

\begin{proof}
A journal of either kind has no mark above the version of the attempt that reads it, because
the versions do not go down. So the current version of attempt $j$ never exceeds $c_j$, and a
gate of $e$ with a version above $c_j$ takes its second branch, which is what
$\strip{c_j}{e}$ has in its place. The two attempts differ in one respect: at such a gate, $e$
brings its version up to date and may write the mark $\ver{c_j}$, where $\strip{c_j}{e}$
writes this mark at its next gate or request, or not at all if the attempt ends or dies
first. The mark stands before the same request in both journals, or it is redundant. A
redundant mark does not change what a later attempt reads, since the next mark or the end
follows it. Typing is preserved, since the second branch of a gate is typed in the same
context as the gate. The last claim follows from Theorem~\ref{thm:stuck} for $e$.
\end{proof}

So all versions of a program are one text, and a deployment is a change of a number. Both
conditions of the theorem are needed. A change that is not behind a gate gets a run stuck
that has passed it, as in Section~\ref{sec:prompt}. And if the versions go down, the old text
meets requests behind a mark that it has no gate for. The artifact has an example of each.

\section{Consequences}
\label{sec:laws}

\subsection{A Long Task Needs a Journal}

\begin{table}[t]
\caption{The expected number of model calls for a task of $n$ calls in a row, if a process
dies during each call with probability $0.01$.}
\label{tab:calls}
\small
\begin{tabular}{@{}lrrr@{}}
\toprule
 & $n = 10$ & $n = 100$ & $n = 1{,}000$ \\
\midrule
start again from nothing & \exScratcha & \exScratchb & \exScratchc \\
replay a journal & \exJournala & \exJournalb & \exJournalc \\
\bottomrule
\end{tabular}
\end{table}

Let a task consist of $n$ calls of the model, one after the other, and let the process die
during each call with probability $q$, after the call was made and before its reply is
stored. With a journal, Theorem~\ref{thm:cost} says that a death repeats one call. Each call
is made until its reply is stored, so the expected number of calls is $n / (1 - q)$. Without
a journal, a new process starts from nothing. If $E_i$ is the expected number of calls still
to be made when $i$ replies are in hand, then $E_n = 0$ and
$E_i = 1 + (1 - q)\,E_{i+1} + q\,E_0$, which gives
\[
E_0 = \frac{(1 - q)^{-n} - 1}{q}.
\]
The first grows linearly with the length of the task and the second exponentially.
Table~\ref{tab:calls} shows both for $q = 0.01$. At a hundred calls, starting again costs
\exScratchb{} calls where the journal costs \exJournalb{}. At a thousand calls it costs more
than two million. An agent that is meant to work for a day cannot do without a journal, and
then it has to be a program for which the journal fits.

Replay itself calls no model, but it rebuilds the context of the agent, and the next call
sends that context to the provider. A provider's cache of prompts is looked up by the text
that is sent~\citep{anthropic2026}. Replay reproduces the earlier requests exactly, so a
process that restarts while the cache still holds the context reads it at the price of a
cached read; a later one pays the full price once~\citep{tpp6}.

\subsection{Time and Chance Belong in the Journal}

A program has two ways to read the clock. $\kw{raw}\;\nm{clock}$ gives the time of the
process that happens to run, and $\kw{do}\;\nm{clock}$ gives the time at which the run first
came to this point. Only the second may reach a request. The same holds for a random number,
for an identifier that is made up on the spot, and for a call of a model. A model need not
give the same answer twice, so an answer that is not journaled is volatile, and a prompt that
is built from it is a volatile argument.

The types do not forbid volatile values. They may be logged, shown, and returned. What the
types say is where the line runs: on one side the values that the journal fixes, on the other
side everything else, with no way from the second to the first.

\subsection{A Prompt Is Code, and a Deployment Goes One Way}

By Theorem~\ref{thm:versions}, an edit is safe if it stands behind a gate, and the edit of a
prompt is no exception. A gate also says when the edit takes effect: for each run, at the
point that the run had reached when a process of the new version first went live. A run never
changes its past, and every attempt replays that past from its start. It follows that the
old branch of a gate can be deleted only when every run that began under an older version has
ended.

A deployment cannot be taken back for the runs that have passed its mark. For them, taking it
back is a new change, and it needs a gate with a higher version whose new branch holds the old
behavior.

\subsection{After a Lost Reply}

The second branch of a $\kw{once}$ has three honest things to do. It can look, if the world
can be asked whether the action happened. It can ask a person; a companion paper types the
question and the answer~\citep{tpp7}. Or it can give up and say so. What it cannot do is try
again as if nothing had happened, since the type of $\kw{do}$ excludes the action. If the
action was approved by a person and the approval is a value that is used up by acting, as
in~\citet{tpp4}, the lost reply leaves the approval used: a second payment needs a second
answer.

A retry in the sense of~\citet{tpp3}, where a generator is called again until a checker
accepts, is a different matter. Each call is a $\kw{do}$ of class $\rd$, and the journal
holds every attempt, so that a new process does not generate again what was already rejected.

\section{Testing the Claims}
\label{sec:tests}

The artifact is one Python script without dependencies. It contains the checker of
Figure~\ref{fig:typing}, the machine of Section~\ref{sec:machine} with a hook at every point
at which a process may die, the reference semantics of Figure~\ref{fig:reference},
Definition~\ref{def:trace} as a function, a generator of random programs, and the tests. All
numbers are those of one fixed seed, and \texttt{make -C artifact check} computes them again
and compares. The script also checks Table~\ref{tab:verdicts} and the three runs of
Section~\ref{sec:overview}, entry by entry.

\subsection{Programs and Runs}

The generator builds programs by following the typing rules. It keeps the labels of the
variables and of the tests above, and chooses among the forms that these allow; a program has
a nesting depth of at most six, and the median program has \exMedianNodes{} nodes in its
syntax tree. A switch makes the generator \emph{slip}: with a probability of one in ten, at
each place where a rule constrains a choice, it ignores the constraint.

A run is random in three ways. The replies of the world are uniform over four numbers. A
process dies with probability $0.15$ before each write and each journaled action, and with
probability $0.02$ where it only reads, up to four times in a run. And each new attempt has
the version of its predecessor or, with probability $0.3$, the next one, up to version $3$.

\subsection{The Theorems}

We took $2{,}000$ programs from the generator without slips, and $2{,}000$ programs from the
generator with slips that the checker accepts. Each of the \exPrograms{} programs was run five
times. The tests are these.
\begin{itemize}
\item Theorem~\ref{thm:stuck}: no run is stuck, and every journal is well formed. A sixth run
  of each program has versions that go up and down, and is not stuck either.
\item Theorem~\ref{thm:sound}: the script reads the answers and the marks from the final
  journal, runs the reference semantics with them and with the $\kw{raw}$ replies of the last
  attempt, and checks that it makes the requests of the journal, ends with the value of the
  run, and has a trace to which the world trace belongs; that no action of class
  $\onceclass$ is repeated; and that no effect lies outside the journal.
\item Theorem~\ref{thm:cost}: repeated performances plus lost replies are at most the
  processes that died.
\end{itemize}
In the \exRuns{} runs, \exCrashes{} processes died, and the journals hold \exRequests{}
requests and \exMarks{} marks. There were \exRepeats{} repeated performances and \exLost{}
lost replies. No test failed.

For Theorem~\ref{thm:complete} the script makes up a reference execution for each of
\exReferences{} programs, with random answers, of which \exRefLost{} are lost, and a random
schedule. It then drives the machine as the proof says and checks that the run has the same
value, the same answers and versions in its journal, and a world trace without repetition.
For Theorem~\ref{thm:versions} it runs $2{,}000$ programs twice on the same random stream,
once with the text of each attempt's version and once with the newest text, and compares
values, world traces, and journals without redundant marks.

\subsection{Every Rule and Every Step Is Needed}
\label{sec:needed}

\begin{table}[t]
\caption{Ten checkers with one weakened rule and five machines with one weakened step. For
each there is a program, accepted by the checker at hand, with a run that the theorems
exclude. The last column is the number of random programs that were tried until the first
such program.}
\label{tab:broken}
\small
\begin{tabular}{@{}lllr@{}}
\toprule
 & What is allowed & What goes wrong & Tried \\
\midrule
\textsc{Do} & a volatile argument & stuck & \exTryDoArg \\
\textsc{Do} & a request under a volatile test & stuck & \exTryDoPc \\
\textsc{Do} & an action of class $\onceclass$ & an action twice & \exTryDoClass \\
\textsc{Once} & a volatile argument & stuck & \exTryOnceArg \\
\textsc{Once} & a request under a volatile test & stuck & \exTryOncePc \\
\textsc{Raw} & a durable label for the reply & stuck & \exTryRawLabel \\
\textsc{Raw} & an action with an effect & an effect per attempt & \exTryRawClass \\
\textsc{If} & branches without the raised $pc$ & stuck & \exTryIfPc \\
\textsc{If} & a value without the label of its test & stuck & \exTryIfResult \\
\textsc{Iter} & a state with the label of its start & stuck & \exTryIter \\
\midrule
row 1 & the request is written after the action & an action twice & \exStepLateReq \\
row 5 & an open $\kw{once}$ is performed again & an action twice & \exStepOnceRetry \\
marks & a gate reads the deployed version & stuck & \exStepNoEpoch \\
row 6 & a replayed request is not compared & a wrong account & \exStepLenient \\
the end & an attempt may end with requests left over & a wrong account & \exStepNoLeftover \\
\bottomrule
\end{tabular}
\end{table}

Table~\ref{tab:broken} lists ten weakenings of Figure~\ref{fig:typing}. For each the script
draws programs from the generator with slips, here with a probability of $0.15$, until the
weakened checker accepts one and one of four runs of it is stuck or violates the second half
of Theorem~\ref{thm:sound}. It finds such a program for every weakening. The slowest needs
\exSlowest{} programs. It is the rule that the value of a conditional carries the label of
its test: the counterexample has to bind such a value, pass it to a request, die after that
request, and read another time in the next process.

The lower part of the table weakens the machine. The first three are tried on programs that
the unweakened checker accepts, and they break its guarantees: an action of class
$\onceclass$ is performed twice, or a run is stuck after a deployment. The last two are tried
on arbitrary programs and break the first half of Theorem~\ref{thm:sound}: a run ends, and no
reference execution has its world trace. This is what the strict comparison in row 6 and at
the end of an attempt is for.

\subsection{Crashes in Tests Find These Defects Rarely}
\label{sec:late}

A team without a type checker would test its agent by killing processes. How much would that
find? We took $2{,}000$ programs from the generator with slips that the checker rejects, and
ran each twenty times with random crashes, as above.

Of the \exRejectedRuns{} runs, \exStuck{} were stuck, and \exAstray{} ended with an effect
too many: an action of class $\onceclass$ performed twice, or an effect outside the journal
repeated by a later process. Together that is \exBadPercent{} percent of the runs. Of the
runs that were stuck, \exLate{} had already performed an action with an effect, which is
\exLatePercent{} percent; such a run can neither go on nor be undone. And \exQuiet{} of the
$2{,}000$ programs, \exQuietPercent{} percent, showed nothing in twenty runs. Such a program
either has a defect that needs a death at one particular step together with a changed
volatile value, or it is safe on every run for a reason that the rules do not see. We did not
separate the two cases. Every run that ended, of every program, satisfied the first half of
Theorem~\ref{thm:sound}.

\subsection{The Cost of a Long Task}

The script checks the closed form of Section~\ref{sec:laws} against its recurrence with exact
fractions, for six values of $n$. It also runs a chain of twenty model calls on the machine,
with a death after one call in ten: over $4{,}000$ runs the chain makes \exChainCalls{} calls
on average, and $20 / 0.9$ is $22.22$.

\subsection{What the Tests Do Not Show}

The tests run small programs over four numbers and a handful of runs of each. They cannot
find an error that needs a longer program or a rarer coincidence. They share the reading of
the definitions with the proofs, since both come from the same author. And they say nothing
about a real engine: the machine of the script is the machine of this paper.

\section{Limits}
\label{sec:limits}

\paragraph{One thing at a time.} The calculus is sequential. Agents make several tool calls
at once and wait for the first reply, and then the order of the answers is itself something
that the journal must fix. We expect the approach to carry over, with the order of arrival as
one more durable value, and have not done it.

\paragraph{The journal grows.} Every attempt replays the whole journal, and a long run has a
long journal. Engines cut it by storing a state and starting a new journal from there. That is
a change of the program text in the middle of a run, and it is outside what
Theorem~\ref{thm:versions} covers.

\paragraph{Classes are claims.} That an action is idempotent is a fact about the world, and
the checker takes it on trust. A wrong class gives a wrong guarantee. The classes are also
coarse: the theorems promise that repetitions are consecutive among the actions of this run,
and say nothing about what other parties do to the world in between.

\paragraph{The model of failure.} A process stops and is gone. The journal is written entry
by entry, and an entry is there or not. A journal that loses or reorders entries, or two
processes that run the same attempt at once, are outside the model.

\paragraph{The language.} Programs are first order and end. An agent that loops without a
bound would need the theorems for the finite parts of a run, which we expect to hold. And a
real agent is written in a real language, where the volatile sources are many: the clock, the
environment, the order in which a hash table is traversed. A checker for such a language has
to find them all; this paper says what it then has to enforce.

\paragraph{Proofs.} The proofs are on paper, and the artifact tests them and does not prove
them.

\section{Related Work}
\label{sec:related}

\paragraph{Durable execution.} Workflow engines replay a history of events against code that
must be deterministic~\citep{temporal-workflow,azure-constraints}, and they offer patching
and versioning for code that changes while runs are
alive~\citep{temporal-patching,azure-versioning}. The constraints are stated in
documentation and enforced at run time by the comparison that we model in row 6.
\citet{burckhardt2021} give a formal model of durable functions without faults, and prove
that execution models with separate storage and with record and replay are equivalent to it.
Our reference semantics has the role of their fault-free model. What we add is a type system for the determinism that
replay presupposes, the two forms of request with the lost reply as a value that the program
handles, and gates.

\paragraph{Idempotence.} \citet{ramalingam2013} show how to make a program idempotent, so
that running it again after a failure is harmless, and prove that failures then cannot be
observed; \citet{jangda2019} give a semantics of serverless platforms in which functions are
run at least once, and conditions under which this equals a simple semantics;
\citet{kallas2023} give a framework with transactions whose implementations on such
platforms are proved correct when requests are executed again. These works make repetition
harmless. We keep the action that must not be repeated in the language, and give
its failure a branch. That a failure between an action and its record cannot be hidden without
the help of the application is the end-to-end argument~\citep{saltzer1984}.

\paragraph{Recovery by logging.} Rollback recovery by message logging relies on the
assumption that a process is deterministic between the events that are
logged~\citep{elnozahy2002}. Lemma~\ref{lem:durable} is this assumption as a property that a
type system establishes. The same survey describes the problem of output that cannot be
taken back, which is our class $\onceclass$. A long-lived transaction that cannot hold locks
can be split into steps with compensations~\citep{garciamolina1987}; the second branch of a
$\kw{once}$ is where a compensation would go.

\paragraph{Information flow.} The type system is that of \citet{volpano1996} with two labels,
in the tradition surveyed by \citet{sabelfeld2003}. The usual reading is secrecy or
integrity. Here the low level is what the journal fixes, and noninterference means that
replay reproduces the requests. A companion paper uses the same technique for integrity,
against text that an attacker has written~\citep{tpp7}.

\paragraph{Updating programs that run.} Dynamic software updating replaces code in a running
program at points where the old state can be carried over~\citep{hicks2005}, and
\citet{stoyle2005} give a calculus in which updates at such points preserve typing. In our
setting the state is not carried over. It is recomputed from the journal, so the new code must
be able to play the old part, and a gate says how. The point of the update is a place in the
journal.

\paragraph{Agents.} The cost of an agent loop, with its cache, is the subject
of~\citet{tpp6}; the journal is what keeps that cost linear when processes die. Consent as a
value that is used up~\citep{tpp4} meets the lost reply in Section~\ref{sec:laws}.

\section{Conclusion}

A journal turns a run of many processes into one run. For this to work, the program must ask
the same questions when it is run again, and it must say for each action whether a second
performance is harmless. We have given both obligations a place in the language: two labels
for the first, two forms of request for the second, and a gate for the case that the program
changes on the way. In return the many ways in which processes can die reduce to one event
that the program has to expect in any case, a reply that does not arrive.

\bibliographystyle{ACM-Reference-Format}
\bibliography{refs}

\end{document}
